\section*{الفصل \ref{ch:g7:sound-production} --- الصوت: إنتاجه وانتشاره}

\begin{solution}{exo:g7:sound-production:1}
المادة التي يسافر فيها \omterm{def:g3:sound-around-us:sound}{الصوت} --- \omterm{def:g2:air-around-us:air}{الهواء} والماء والخشب
والفولاذ، وأيّ حشد جزيئي. واللاوسط: الخلاء.
\end{solution}

\begin{solution}{exo:g7:sound-production:2}
التناقل هو الذي يسافر: نمط منطلق من التضاغطات يُسلَّم من صفّ إلى
صفّ في الحشد. وكل \omterm{def:g7:states-of-matter:molecule}{جزيء} لا يفعل إلا أن يهتزّ حول موضعه ---
فالرسالة تتحرّك، والرُّسل يمكثون.
\end{solution}

\begin{solution}{exo:g7:sound-production:3}
لا حشد، فلا تناقل: فالجرّة المشفوطة تترك جدران الجرس المرتعشة
تدفع في لا شيء، فلا تتكوّن رسالة. ومعارك الفضاء صامتة للسبب
نفسه --- فالانفجارات بين النجوم لا تجد حشدًا يحمل هديرها.
\end{solution}

\begin{solution}{exo:g7:sound-production:4}
\omterm{def:g2:air-around-us:air}{الهواء} نحو \qty{340}{m/s}؛ والماء نحو \qty{1500}{m/s}؛
والفولاذ نحو \qty{5000}{m/s}. وكلما اشتدّ تلامس \omterm{def:g7:states-of-matter:molecule}{الجزيئات} أسرع
تسليم الدفعة: طائرة، فمتلامسة، فمقفلة.
\end{solution}

\begin{solution}{exo:g7:sound-production:5}
صفوف الفولاذ المقفلة تناقل يدًا بيد تقريبًا: نحو خمسة عشر ضعف
وتيرة \omterm{def:g2:air-around-us:air}{الهواء}، فتصل رسالة القضيب قبل الهدير المحمول في \omterm{def:g2:air-around-us:air}{الهواء}
بزمن طويل.
\end{solution}

\begin{solution}{exo:g7:sound-production:6}
رحلة \omterm{def:g1:light-and-shadows:source}{الضوء} عبر أيّ ملعب تُعدّ لحظية بجانب رحلة \omterm{def:g3:sound-around-us:sound}{الصوت}: فالضربة
المرئية تعلّم البداية الحقيقية بدقّة أعلى بكثير مما تستطيع يد
ساعة الإيقاف أن تستعمله.
\end{solution}

\begin{solution}{exo:g7:sound-production:7}
$340 \times 1 = \qty{340}{m}$ ذهابًا وإيابًا: فالجدار يقف على
\qty{170}{m}. والجواب \qty{340}{m} ينسى أن الثانية دفعت ثمن
الذهاب \emph{والعودة} معًا.
\end{solution}

\begin{solution}{exo:g7:sound-production:8}
في الماء: $1500 \times 2 = \qty{3000}{m}$ ذهابًا وإيابًا ---
فالعمق \qty{1500}{m}. والزلّة: استعارة \qty{340}{m/s} الهوائية
لرحلة تحت الماء.
\end{solution}

\begin{solution}{exo:g7:sound-production:9}
$d = 340 \times 2.5 = \qty{850}{m}$. فوميض اللعبة النارية أعطى
مسدّس انطلاق حقيقيًّا؛ أما الطائرة غير المرئية فلا تعطي شيئًا
--- فلا لحظة إصدار يُوقَّت منها --- والأسوأ أن الهدير المسموع
غادر الطائرة منذ ثوانٍ: فالأذن تشير إلى حيث \emph{كانت}
الطائرة.
\end{solution}

\begin{solution}{exo:g7:sound-production:10}
الأذن المغمورة يخدمها تناقل الماء \omterm{def:g7:motion-average-speed:formula}{بسرعة} \qty{1500}{m/s}، وأذن
الحافة يخدمها تناقل \omterm{def:g2:air-around-us:air}{الهواء} \omterm{def:g7:motion-average-speed:formula}{بسرعة} \qty{340}{m/s} --- وكثير من
\omterm{def:g3:sound-around-us:sound}{صوت} المكبّر لا يعبر حدّ السطح إلى \omterm{def:g2:air-around-us:air}{الهواء} البتّة. \omterm{def:g7:sound-production:medium}{وسط} أسرع،
ورسالة خاصّة: فيفوز السابح رغم المسافة.
\end{solution}

\begin{solution}{exo:g7:sound-production:11}
ضربة واحدة، وتناقلان: عبر الفولاذ وعبر \omterm{def:g2:air-around-us:air}{الهواء}. الفولاذ:
$3000 \div 5000 = \qty{0.6}{s}$؛ \omterm{def:g2:air-around-us:air}{والهواء}: $3000 \div 340 \approx
\qty{8.8}{s}$ --- فالرنين يتقدّم الدويّ بنحو \qty{8.2}{s}.
\end{solution}

\begin{solution}{exo:g7:sound-production:12}
قرعة واحدة ترسل اللحظة نفسها في الوسطين معًا؛ وعند \qty{750}{m}
يلتقط المسجّل وصول الماء (على اللاقط المائي) ووصول \omterm{def:g2:air-around-us:air}{الهواء} (على
المِذياع). والأزمنة المتوقّعة: الماء $750 \div 1500 =
\qty{0.5}{s}$، \omterm{def:g2:air-around-us:air}{والهواء} $750 \div 340 \approx \qty{2.2}{s}$ ---
أي فارق نحو \qty{1.7}{s}. وبالعكس: إذ تُعرف $d$ \omterm{def:g7:motion-average-speed:formula}{وسرعة} \omterm{def:g2:air-around-us:air}{الهواء}،
يعطي الفارق المقيس زمن وصول الماء ($t_{\text{الهواء}} -
\text{الفارق}$)، فيكون $v = 750 \div 0.5 = \qty{1500}{m/s}$ ---
وتيرة الماء مقيسة من غير توقيت القرعة الصامتة نفسها البتّة.
\end{solution}

\begin{solution}{pb:g7:sound-production:1}
\textbf{1.} $340 \times 4 = \qty{1360}{m}$ ذهابًا وإيابًا:
فالعرض \qty{680}{m}.
\textbf{2.} $340 \times 1.5 = 510$؛ والنصف: \qty{255}{m}.
\textbf{3.} الجدار الأول: $340 \times 1 \div 2 = \qty{170}{m}$؛
والثاني: $340 \times 2 \div 2 = \qty{340}{m}$؛ فيمتدّ الأخدود
$170 + 340 = \qty{510}{m}$ على ذلك الخطّ.
\textbf{4.} الصيغة مضبوطة؛ واليد ليست كذلك: فبدء الإنسان
وإيقافه يتبعثران بأعشار الثانية --- أي عشرات الأمتار عند
\qty{340}{m/s}. ومن هنا القاعدة القديمة: كرّر، وراقب القراءات
تتجمّع، وأبلغ عن التجمّع.
\textbf{5.} $1500 \times 0.2 = \qty{300}{m}$ ذهابًا وإيابًا:
فالعمق \qty{150}{m}.
\textbf{6.} $1500 \times 0.4 \div 2 = \qty{300}{m}$.
\textbf{7.} عند حدّ \omterm{def:g2:air-around-us:air}{الهواء} والماء يُردّ معظم الصرخة --- فيدخل
النهر القليل، ولا بدّ لصدى القاع الضعيف أن يعبر الحدّ ثانية كي
يبلغ أذنًا في \omterm{def:g2:air-around-us:air}{الهواء}: فلا ينجو شيء يُذكر. وحتى لو سلّمنا
\omterm{ex:g4:how-sound-travels:echo}{بالصدى}، فرحلات ذهاب وإياب بأعشار الثانية تهزم أيّ توقيت
بالأذن.
\textbf{8.} زمن \omterm{ex:g4:how-sound-travels:echo}{الصدى} يطول فجأة بينما يعبر القارب الهبوط ---
قاع أعمق، فرحلة أطول: فالسونار يرسم الجُرف تحت الماء.
\textbf{9.} الفولاذ: $2000 \div 5000 = \qty{0.4}{s}$؛ \omterm{def:g2:air-around-us:air}{والهواء}:
$2000 \div 340 \approx \qty{5.9}{s}$.
\textbf{10.} التناقلان يبدآن معًا، فينمو الفارق بالتناسب مع
المسافة: فكل كيلومتر يضيف تأخيرًا إضافيًّا ثابتًا للتناقل
البطيء على السريع. فقِس الفارق، واقسمه على الفرق لكل كيلومتر،
تسقط المسافة --- ضربة مطرقة واحدة، بلا حاجة إلى ساعة إيقاف عند
الطرف البعيد.
\textbf{11.} \omterm{def:g3:sound-around-us:sound}{صوت} القضيب: فالضباب حشد من القطيرات يبعثر
\emph{\omterm{def:g1:light-and-shadows:source}{الضوء}} بلا أمل، لكنه لا يكاد يزعج تناقلًا مقفلًا داخل
الفولاذ الصلب (\omterm{def:g3:sound-around-us:sound}{وصوت} \omterm{def:g2:air-around-us:air}{الهواء} أيضًا ينجو من الضباب نجاة حسنة ---
لكن رسالة القضيب أنظف الثلاث وأسرعها).
\textbf{12.} مثلًا: ``كل صدى يدفع ثمن الذهاب والإياب: فنصِّف
قبل أن تصدّق. وكل \omterm{def:g7:sound-production:medium}{وسط} يضع وتيرته: فاختر \omterm{def:g7:motion-average-speed:formula}{السرعة} مع الحشد، لا
بالعادة. وكل توقيت في هذا الأخدود بدأ بالعين --- فرحلة \omterm{def:g1:light-and-shadows:source}{الضوء}
تُعدّ لحظية هنا، ولذلك تصلح رؤية المطرقة تهوي مسدّسَ
انطلاق.''
\end{solution}
